% ============================================================================
% Methods (from Arce's gBKMR_Nature.docx)
% NOTE: Formulas were stripped during docx→text conversion. TODO markers
% below indicate places to paste the LaTeX formulas from Zilan's original
% .tex file (statistics in medicine draft / dissertation).
% ============================================================================

\section*{Methods}

\subsection*{Identifiability assumptions and estimation}

We consider the setting in which $n$ subjects, labelled $i = 1, \ldots, n$, enter a study at baseline (time $t = 0$) and are exposed to $A_{t,i} = A_{0,i}$. In a clinical trial, the value of $A_{0,i}$ is determined by the study protocol, usually at random. In an observational study, $A_{0,i}$ is observed and recorded. A collection of covariates, $C_{0,i}$, are also measured at baseline. At the subsequent follow-up visits, the time-varying confounders $L_{t,i}$ and treatment $A_{t,i}$ are recorded, $t = 1, \ldots, T$. An outcome $Y_i$ is observed for each subject, measured at the final visit $T$. $Y_i$ can be continuous or binary.

Along with the observed outcome $Y_i$ for each subject, for each possible realization of the treatment history $\mathbf{a} \in \mathcal{A}$, we also define $Y_i^{\mathbf{a}}$, the potential outcome that would have been observed had the subject, possibly contrary to fact, received treatment history $\mathbf{a}$. In a simple setting with 3 visits, we assume the Directed Acyclic Graph (DAG) in Figure~\ref{fig:dag}. Under a static, fixed treatment regime, the overall causal effect per one interquartile range increase in all exposures (i.e.\ changing from the 25th percentile to the 75th percentile) can be defined as:

% TODO: paste ACE equation from Zilan's .tex
% ACE_{a,a*} = E[Y^{A_0=a_0*, A_1=a_1*, A_2=a_2*}] - E[Y^{A_0=a_0, A_1=a1, A_2=a_2}]

Where $a_0, a_1, a_2$ are the 25th percentile of the exposures, and $a_0^*, a_1^*, a_2^*$ are the 75th percentile of the exposures. In addition to estimating the effect of the mixture per one interquartile range increase (or per any comparison we choose), we could also estimate the effect of a single exposure of the mixture across time points.

The basic idea of the g-formula is standardization, a common approach to estimating the expected value of $Y^{\mathbf{a}}$ in the population. It can be illustrated in a simple time-fixed confounding example with baseline confounders $C_0$ where the change in the observed mean of the outcome under two different exposure levels is standardized over the observed values $c_0$ of $C_0$.

% TODO: paste equation E[Y^a] = sum_{c0} E[Y|A=a, C_0=c_0] P(C_0=c_0)

where the sum is over all possible values $c_0$ of $C_0$. Following the DAG in Figure~\ref{fig:dag}, the following identifiability assumptions are required for the g-formula to be a valid estimator for the average causal effect:

\paragraph{Consistency:}
% TODO: paste consistency equation

\paragraph{Exchangeability:}
% TODO: paste sequential exchangeability equations (3 conditions)

\paragraph{Positivity:}
% TODO: paste positivity equation

The general expression of the g-formula for a static fixed intervention when the identifiability assumptions hold is:

% TODO: paste full g-formula expression

\subsection*{Bayesian Kernel Machine Regression}

We first review Kernel Machine Regression (KMR) as a framework for estimating the effect of a complex mixture when only a single time point of exposure is measured. For each subject $i = 1, \ldots, n$, we assume:

% TODO: paste Y_i = h(a_i) + c_i^T beta + epsilon_i

where $a_i = (a_{i1}, \ldots, a_{iM})^T$ is a vector of $M$ exposure variables. Liu et al.\ showed that this model can be expressed as the mixed model:

% TODO: paste y_i ~ N(h_i + c_i^T beta, sigma^2) and h ~ N(0, tau K)

where $K$ is the kernel matrix, which has $(i, j)$ element $K(a_i, a_j)$.

We focus on the Gaussian kernel, which flexibly captures a wide range of underlying functional forms for $h(\cdot)$, although the methods are applicable to a broad choice of kernels. Under the previous model, we assume:

% TODO: paste corr(h_i, h_j) = exp(-1/rho * sum_m (a_{im} - a_{jm})^2)

Where $\rho$ is a tuning parameter that regulates the smoothness of the dose-response function. This assumption implies that two subjects with similar exposures will have more similar risks (i.e.\ if $a_i$ is close to $a_j$, then $h_i$ will be close to $h_j$).

Bobb et al.\cite{bobb_bkmr_2015} presented the Bayesian Kernel Machine Regression (BKMR) as a framework to estimate the effect of a complex mixture on a health outcome. To fit BKMR, we assume a prior on the coefficients for the confounding variables, $\beta \sim 1$ (flat prior), and $\sigma^{-2} \sim \text{Gamma}(a_\sigma, b_\sigma)$, where we set $a_\sigma = b_\sigma = 0.001$. We can parameterize BKMR by $\lambda = \tau \sigma^{-2}$, and we assume a Gamma prior distribution for the variance component of $\lambda$. For the smoothness parameter $\rho$, we assume $\rho \sim \text{Unif}(a, b)$ with $a = 0$ and $b = 100$. For additional details regarding BKMR and prior specifications, see Bobb et al.\cite{bobb_bkmr_2015}. Note that the confounding variables can be included in the $h(\cdot)$ function to allow more flexibility in the confounding adjustment.

Utilizing the Bayesian model, we are able to calculate the posterior inclusion probability (PIP) of the exposure variables, which denotes the likelihood of a particular exposure being included in the model. % TODO: ADD VARIABLE SELECTION description (spike-and-slab)

\subsection*{g-BKMR}

g-BKMR is a Bayesian semi-parametric regression method that allows for (i) estimation of joint effects and interactions of time-varying environmental exposures, (ii) time-varying confounding. For the g-formula to be implemented, we must estimate $E[Y \mid A, L, C_0]$ and $f(L_t \mid \bar{A}_{t-1}, \bar{L}_{t-1}) f(C_0)$ from the observed data. This involves postulating a model for $Y$ given $A$ and $L$ and models for $L_t$ given $\bar{A}_{t-1}$ and $\bar{L}_{t-1}$. Using g-BKMR, we can estimate the causal effect for a change in exposure profile from $\mathbf{a}$ to $\mathbf{a}^*$ via the following algorithm:

\begin{enumerate}
    \item Fit BKMR time-varying confounding and outcome models:

    % TODO: paste equations (4) and (5) from Zilan's .tex

    Note that multiple time-varying confounders at multiple time-points are allowed. Each of these would be modeled under the BKMR framework and can be fitted in parallel to speed computation.

    \item For each Markov chain Monte Carlo (MCMC) iteration,

    \begin{enumerate}
        \item Generate $K$ samples of the confounder for the mean level of covariates under exposure level $\mathbf{a}$ from confounder model (4).
        \item For each of the MCMC iterations and $K$ samples of $L$, estimate the average outcome value for the mean level of covariates for $\mathbf{a}$ from model (5).
        \item Let the posterior sample of $Y^{\mathbf{a}}$ be the mean over the $K$ samples.
    \end{enumerate}

    \item The potential outcome is then estimated as:
    % TODO: paste estimator expression

    \item The posterior sample of the ATE is obtained by:
    % TODO: paste ATE posterior expression

    \item Estimate the 95\% credible intervals from these posterior samples.
\end{enumerate}

The implementation of the g-BKMR algorithm discussed in this paper is publicly available on GitHub at the following link: \url{https://github.com/zc2326/causalbkmr}. In the current implementation, binary or continuous outcomes and time-varying covariates are accommodated.
